\documentclass[11pt,a4paper]{article}
\usepackage[utf8]{inputenc}
\usepackage{lmodern}
\usepackage[T1]{fontenc}
\usepackage[finnish]{babel}
\usepackage{graphicx}
\usepackage{epstopdf}
\usepackage{listings} 
\usepackage{xcolor}
\usepackage{textcomp}

\definecolor{dkgreen}{rgb}{0,0.6,0}
\definecolor{gray}{rgb}{0.5,0.5,0.5}
\definecolor{amber}{rgb}{1.0, 0.49, 0.0}

%Java lstset
\lstset{frame=none,
  language=Java,
  aboveskip=3mm,
  belowskip=3mm,
  showstringspaces=false,
  columns=flexible,
  basicstyle={\small\ttfamily},
  numbers=none,
  numberstyle=\tiny\color{gray},
  keywordstyle=\color{blue},
  commentstyle=\color{dkgreen},
  stringstyle=\color{amber},
  breaklines=true,
  breakatwhitespace=true,
  tabsize=3
}

%Multi-byte encoding
\lstset{literate=
  {á}{{\'a}}1 {é}{{\'e}}1 {í}{{\'i}}1 {ó}{{\'o}}1 {ú}{{\'u}}1
  {Á}{{\'A}}1 {É}{{\'E}}1 {Í}{{\'I}}1 {Ó}{{\'O}}1 {Ú}{{\'U}}1
  {à}{{\`a}}1 {è}{{\`e}}1 {ì}{{\`i}}1 {ò}{{\`o}}1 {ù}{{\`u}}1
  {À}{{\`A}}1 {È}{{\'E}}1 {Ì}{{\`I}}1 {Ò}{{\`O}}1 {Ù}{{\`U}}1
  {ä}{{\"a}}1 {ë}{{\"e}}1 {ï}{{\"i}}1 {ö}{{\"o}}1 {ü}{{\"u}}1
  {Ä}{{\"A}}1 {Ë}{{\"E}}1 {Ï}{{\"I}}1 {Ö}{{\"O}}1 {Ü}{{\"U}}1
  {â}{{\^a}}1 {ê}{{\^e}}1 {î}{{\^i}}1 {ô}{{\^o}}1 {û}{{\^u}}1
  {Â}{{\^A}}1 {Ê}{{\^E}}1 {Î}{{\^I}}1 {Ô}{{\^O}}1 {Û}{{\^U}}1
  {Ã}{{\~A}}1 {ã}{{\~a}}1 {Õ}{{\~O}}1 {õ}{{\~o}}1
  {œ}{{\oe}}1 {Œ}{{\OE}}1 {æ}{{\ae}}1 {Æ}{{\AE}}1 {ß}{{\ss}}1
  {ű}{{\H{u}}}1 {Ű}{{\H{U}}}1 {ő}{{\H{o}}}1 {Ő}{{\H{O}}}1
  {ç}{{\c c}}1 {Ç}{{\c C}}1 {ø}{{\o}}1 {å}{{\r a}}1 {Å}{{\r A}}1
  {€}{{\euro}}1 {£}{{\pounds}}1 {«}{{\guillemotleft}}1
  {»}{{\guillemotright}}1 {ñ}{{\~n}}1 {Ñ}{{\~N}}1 {¿}{{?`}}1
}

\begin{document}
\title{Tietokantojen perusteet, kevät 2020 \\
Harjoitustyön raportti}
\author{Ville Manninen\\ \\014922187 \\macville}
\date{\today}
\clearpage\maketitle
\begin{center}
\vspace{2cm}
\includegraphics{university-of-helsinki-2.eps} 
\end{center}
\thispagestyle{empty}
\newpage
\tableofcontents
\newpage
\section{Esittely}
\subsection*{Yleistä}
Harjoitustyö on toteutettu Java-ohjelmana käyttäen NetBeans ympäristöä, ohjelmalla on seuraavat ominaisuudet.
\begin{itemize}
\item Tietokannan ja puuttuvien taulukoiden luominen.
\item Asiakkaiden lisääminen tietokantaan nimellä.
\item Paikan lisäämien tietokantaan nimellä.
\item Paketin lisääminen tietokantaan käyttäen asiakkaan nimeä.
\item Tapahtuman lisääminen tietokantaan käyttäen paikannimeä ja paketin seurantakoodia.
\item Paketin tapahtumien hakeminen
\item Asiakkaan pakettien hakeminen ja yksittäisten pakettien tapahtumien lukumäärä'
\item Tapahtumien määrän hakeminen paikasta tietyllä päivämäärällä
\end{itemize}
\newpage
\subsection*{Komennot}
Ohjelmalle annetaan komentoja komentoriviltä käyttäen pelkästään komennon numeroa, komennon jälkeen ohjelma kysyy käyttäjältä tarvittavia jatkosyötteitä. Ohjelmalla on seuraavat komennot jotka esiintyvät myös ohjelman valikossa.
\begin{description}
\item \textbf{1} Uuden asiakkaan lisääminen
\item \textbf{2} Uuden paikan lisääminen
\item \textbf{3} Paketin lisääminen
\item \textbf{4} Uuden tapahtuman lisääminen
\item \textbf{5} Tulostaa paketin tapahtumat
\item \textbf{6} Tulostaa asiakkaan paketit ja niihin liittyvien tapahtumien määrän
\item \textbf{7} Tapahtumien määrän haku paikasta tietyllä päivämäärällä
\item \textbf{0} Sulkee ohjelman 
\end{description}

\newpage
\section{Kaaviot ja Skeemat}
\subsection*{Tietokantakaavio}

\subsection*{SQL-skeema}

\newpage
\section{Tehokkuustesti}
\subsection*{Ilman indeksejä}
\subsection*{Indeksien lisäämisen jälkeen}

\newpage
\section{Tietokannan eheys}
Ohjelman tietokannan eheyttä pidetään yllä käyttämällä SQLiten sekä Javan tarjoamia ominaisuuksia.
Tietokantaan ei ole mahdollista luoda samannimisiä asiakkaita tai paikkoja, eikä tietokantaan luoduille paketeille saa antaa samaa seurantakoodia. Samannimisiä syötteitä valvotaan SQLiten rajoitteella UNIQUE, joka varmistaa ettei samaan taulukkoon lisätä samannimisiä arvoja. Lisäksi SQLiten rajoitteella NOT NULL valvotaan että sarakkeisiin ei lisätä tyhjiä arvoja.
 
Java-ohjelma tarkistaa käyttäjän antamat syötteet ja ilmoittaa syyn jos annettu syöte on puutteellinen, esimerkiksi jos käyttäjän syöte on tyhjä. Tietokantaan lisättäessä paketteja tai tapahtumia jotka viittaavat toisiin taulukoihin, sovellus varmistaa että annetut viitteet löytyvät tietokannasta.

\newpage
\section{Lähdekoodi}
%Tarkista lstinputlisting polut
\textbf{https://github.com/Viltska/sql-app-training}.

\subsection*{Main.java}
\lstinputlisting[language=Java]{E:/Koulu/Tietokantojen Perusteet/sql-app-training/src/Main.java}
\subsection*{DatabaseManager.java}
\lstinputlisting[language=Java]{E:/Koulu/Tietokantojen Perusteet/sql-app-training/src/DatabaseManager.java}
\subsection*{Asiakkaat.java}
\lstinputlisting[language=Java]{E:/Koulu/Tietokantojen Perusteet/sql-app-training/src/Asiakkaat.java}
\subsection*{Paikat.java}
\lstinputlisting[language=Java]{E:/Koulu/Tietokantojen Perusteet/sql-app-training/src/Paikat.java}
\subsection*{Paketit.java}
\lstinputlisting[language=Java]{E:/Koulu/Tietokantojen Perusteet/sql-app-training/src/Paketit.java}
\subsection*{Tapahtumat.java}
\lstinputlisting[language=Java]{E:/Koulu/Tietokantojen Perusteet/sql-app-training/src/Tapahtumat.java}
\end{document}